\section{Threats to Validity}\label{sec:threats}

\textbf{Model access.} The models were reached through different routes (Table~\ref{tab:models}), a known risk for studies that use LLMs \citep{sallou2024breaking}. Three were generated through conversational interfaces rather than a raw API: Claude through Claude Code subagents, which add their harness's system prompt to ours, and GPT and DeepSeek through their web chats, whose own system prompts and defaults are outside the study's control. Temperature and the output limit were therefore controlled only for Gemini and Qwen, which affects the consistency comparison of RQ3 more than any other result. GPT was generated with ChatGPT's default model, so its exact version is not known; it answered with a medium reasoning effort that was left unchanged. Its web sessions were operated by a third party following a written protocol (a new temporary, unpersonalised chat per prompt, no tools, no follow-up, replies copied unedited) whose adherence could not be observed directly, and for 16 of its 42 prompts, from the three projects with the longest documentation, the interface turned the pasted prompt into an attached file with identical content. Gemini is a flash-tier model rather than the planned Pro tier, and its hidden reasoning tokens count against its 8,192-token output limit. Qwen's free tier rejects the two largest projects' prompts, so its missing cells reflect access, not quality. A sixth model, Llama, was dropped before the main run because no working access route was found.

\textbf{Training-data contamination.} All projects are public, and only two of the seven manifest repositories have fewer than 2,000 GitHub stars, so a model may reproduce a memorised manifest rather than reason from the documentation. Provider refusals for recitation are recorded, and each model's similarity to the human manifest is compared with its similarity to the other models' configurations for the same project, so that only an excess over the peer level is read as reproduction (Section~\ref{sec:disc-memorisation}). Reproduction would pull a model's configurations toward the human manifest and so shrink the differences reported for RQ1.

\textbf{Selection of projects and of the human baseline.} Results describe these seven projects, not application classes in general (Section~\ref{sec:disc-generalisation}). The set was fixed after a two-project pilot and before the main run. Because a project was selected only if its human manifest deployed on the evaluation cluster, the human D5 outcomes are partly true by construction, and the selection favours applications that run on a local cluster; human runtime results are therefore reported as a reference only.

\textbf{Evaluation cluster.} D5 and D6 are measured on a four-node \texttt{kind} cluster on a single host. Cloud load balancers, multi-zone scheduling and production storage are not exercised, and absolute times should not be read as representative of production. \texttt{kind}'s local-path storage does not apply a pod's \texttt{fsGroup}, which could make a database fail here that would start elsewhere. None of the eight database failures depends on it: seven run the container as root with every capability dropped, which prevents the entrypoint from changing file ownership or switching user whatever the volume's group, and the eighth fails because its data directory is owned by another user, which \texttt{fsGroup} does not change.

\textbf{Measurement adjustments.} The pipeline was corrected 25 times during the study (Appendix~\ref{app:fixes}); each correction was logged with its trigger, applied to every configuration it could affect and followed by re-evaluation. All were made after generation had begun, and because most were found while diagnosing failures, they are one-sided: 8 affected only the running of the pipeline, 3 made scoring stricter, and 14 removed a failure that the investigation traced to the measurement or to an undocumented prerequisite, never a success that turned out to be spurious. The deployment results are robust to leaving out any one project (Section~\ref{sec:res-rq5}). Two adjustments supply facts the bundles lack (image aliases and start-up prerequisites, Section~\ref{sec:method-pipeline}); they apply to 29 generated configurations and to no human manifest. In the worst case, in which every adjusted configuration that succeeded would otherwise have failed, the share of generated configurations both fully Ready and passing the smoke test falls from 59.6\% to 51.0\%. Four of the dashboard failures concern a start-up Secret that the documentation never mentions: once the Secret is supplied, these configurations fail because their RBAC rules do not let the dashboard update it, which a deployer would also have to discover. Three configurations set a required variable to an empty string, which the application rejects; the start-up adjustment fires only when a variable is absent, so they are scored as failing.

\textbf{Labelling and tools.} Root-cause labels were assigned by a single rater, drafted with the help of an AI assistant from the same model family as one of the evaluated models (Claude). The blind second labelling pass was also made by a Claude model, not by a second human: it shows that the labels follow from the recorded evidence ($\kappa=0.84$), but it cannot rule out errors that models of one family share. The evidence level of each label is reported, and every label is published with its evidence so that it can be re-rated. The static tools are imperfect and disagree with one another, so their findings are read as what these pinned tool versions detect; agreement between the two security scanners is reported. Redaction decisions followed a written rule set, with every redaction logged and reviewed by a second author, but borderline cases remain possible. Facts that the bundles do not contain are excluded from D4, while the default ports of widely used database images are treated as general knowledge and scored.

\textbf{Resilience measurement.} The two injected faults are handled by Kubernetes controllers whatever the configuration says, so recovery success barely discriminates, and downtime depends mainly on replica count. D6 therefore characterises the replica choices a model makes rather than a broader notion of resilience; faults such as node drains or resource pressure would test more of a configuration's choices.

\textbf{Drift.} Providers update the weights behind a model identifier without notice, and an unpinned image tag resolves to whatever is current. Generation dates are recorded for every configuration, and images were resolved in November 2026. Tag drift affects the human baselines too: the human Sock Shop manifest references \texttt{mongo} without a tag, which now resolves to a release whose wire protocol the project's 2017 cart service no longer speaks; re-deploying with the image pinned to an older release passed, but the baseline is scored as published.
